
\begin{longtable}[t]{rlrrrrl}
\caption{\label{tab:balance}Covariate balance}\\
\toprule
Study & Outcome & N & F & df1 & df2 & p\\
\midrule
\endfirsthead
\toprule
Study & Outcome & N & F & df1 & df2 & p\\
\midrule
\endhead
1 & party id & 686 & 0.44 & 3 & 682 & 0.727\\
1 & ideology & 685 & 0.30 & 3 & 681 & 0.828\\
1 & education code & 685 & 2.42 & 3 & 681 & 0.065\\
1 & age years & 686 & 1.23 & 3 & 682 & 0.297\\
1 & is female & 686 & 1.53 & 3 & 682 & 0.204\\
1 & is white & 783 & 0.64 & 1 & 781 & 0.424\\
2 & party id & 824 & 0.76 & 3 & 820 & 0.519\\
2 & ideology & 822 & 1.29 & 3 & 818 & 0.275\\
2 & education code & 823 & 0.78 & 3 & 819 & 0.505\\
2 & age years & 824 & 0.56 & 3 & 820 & 0.643\\
2 & is female & 823 & 0.47 & 3 & 819 & 0.706\\
2 & is white & 945 & 1.40 & 1 & 943 & 0.238\\
3 & party id & 1055 & 0.08 & 1 & 1053 & 0.784\\
3 & ideology & 1055 & 0.00 & 1 & 1053 & 0.961\\
3 & education code & 1055 & 2.51 & 1 & 1053 & 0.113\\
3 & age years & 1053 & 1.00 & 1 & 1051 & 0.317\\
3 & is female & 1053 & 1.04 & 1 & 1051 & 0.309\\
3 & is white & 1049 & 0.06 & 1 & 1047 & 0.810\\
\bottomrule
\end{longtable}
\noindent\footnotesize Classical F tests. Studies 1 and 2 use policy-congeniality factorials except race, which uses incentives alone. Study 3 uses incentives alone. Alternative race specifications are in race\_balance\_checks.csv.\par\normalsize
